\documentclass[10pt,a4paper]{amsart}
\usepackage{amsmath,amssymb,amsthm,graphicx}
\usepackage{verbatim}
\newenvironment{reference}{\comment}{\endcomment}
%\newenvironment{reference}{}{}
\newenvironment{slogan}{\comment}{\endcomment}
\newenvironment{history}{\comment}{\endcomment}

% For commutative diagrams we use Xy-pic
\usepackage[all]{xy}

% We use 2cell for 2-commutative diagrams.
\xyoption{2cell}
% Theorem environments.
%
\theoremstyle{plain}
\newtheorem{theorem}[subsection]{Theorem}
\newtheorem{proposition}[subsection]{Proposition}
\newtheorem{lemma}[subsection]{Lemma}

\theoremstyle{definition}
\newtheorem{definition}[subsection]{Definition}
\newtheorem{example}[subsection]{Example}
\newtheorem{exercise}[subsection]{Exercise}
\newtheorem{situation}[subsection]{Situation}

\theoremstyle{remark}
\newtheorem{remark}[subsection]{Remark}
\newtheorem{remarks}[subsection]{Remarks}

\numberwithin{equation}{subsection}

% Macros
%
\def\lim{\mathop{\mathrm{lim}}\nolimits}
\def\colim{\mathop{\mathrm{colim}}\nolimits}
\def\Spec{\mathop{\mathrm{Spec}}}
\def\Hom{\mathop{\mathrm{Hom}}\nolimits}
\def\Ext{\mathop{\mathrm{Ext}}\nolimits}
\def\SheafHom{\mathop{\mathcal{H}\!\mathit{om}}\nolimits}
\def\SheafExt{\mathop{\mathcal{E}\!\mathit{xt}}\nolimits}
\def\Sch{\mathit{Sch}}
\def\Mor{\mathop{\mathrm{Mor}}\nolimits}
\def\Ob{\mathop{\mathrm{Ob}}\nolimits}
\def\Sh{\mathop{\mathit{Sh}}\nolimits}
\def\NL{\mathop{N\!L}\nolimits}
\def\CH{\mathop{\mathrm{CH}}\nolimits}
\def\proetale{{pro\text{-}\acute{e}tale}}
\def\etale{{\acute{e}tale}}
\def\QCoh{\mathit{QCoh}}
\def\Ker{\mathop{\mathrm{Ker}}}
\def\Im{\mathop{\mathrm{Im}}}
\def\Coker{\mathop{\mathrm{Coker}}}
\def\Coim{\mathop{\mathrm{Coim}}}

% Boxtimes
%
\DeclareMathSymbol{\boxtimes}{\mathbin}{AMSa}{"02}

%
% Macros for moduli stacks/spaces
%
\def\QCohstack{\mathcal{QC}\!\mathit{oh}}
\def\Cohstack{\mathcal{C}\!\mathit{oh}}
\def\Spacesstack{\mathcal{S}\!\mathit{paces}}
\def\Quotfunctor{\mathrm{Quot}}
\def\Hilbfunctor{\mathrm{Hilb}}
\def\Curvesstack{\mathcal{C}\!\mathit{urves}}
\def\Polarizedstack{\mathcal{P}\!\mathit{olarized}}
\def\Complexesstack{\mathcal{C}\!\mathit{omplexes}}
% \Pic is the operator that assigns to X its picard group, usage \Pic(X)
% \Picardstack_{X/B} denotes the Picard stack of X over B
% \Picardfunctor_{X/B} denotes the Picard functor of X over B
\def\Pic{\mathop{\mathrm{Pic}}\nolimits}
\def\Picardstack{\mathcal{P}\!\mathit{ic}}
\def\Picardfunctor{\mathrm{Pic}}
\def\Deformationcategory{\mathcal{D}\!\mathit{ef}}
\usepackage{cleveref}
\crefname{originalreference}{original reference}{original references}
\Crefname{originalreference}{Original reference}{Original references}
\makeatletter
\newlabel{lemma-galois-goes-up}{{09I2}{source}{}{sourcepointer-1}{}}
\newlabel{lemma-galois-goes-up@cref}{{[originalreference][0][]09I2}{[1][0][]source}{}{}{}}
\makeatother
\crefname{originalreference}{original reference}{original references}
\Crefname{originalreference}{Original reference}{Original references}
\makeatletter
\newlabel{lemma-galois-infinite}{{0BMK}{source}{}{sourcepointer-2}{}}
\newlabel{lemma-galois-infinite@cref}{{[originalreference][0][]0BMK}{[1][0][]source}{}{}{}}
\makeatother
\crefname{originalreference}{original reference}{original references}
\Crefname{originalreference}{Original reference}{Original references}
\makeatletter
\newlabel{lemma-galois-profinite}{{0BMJ}{source}{}{sourcepointer-3}{}}
\newlabel{lemma-galois-profinite@cref}{{[originalreference][0][]0BMJ}{[1][0][]source}{}{}{}}
\makeatother
\crefname{originalreference}{original reference}{original references}
\Crefname{originalreference}{Original reference}{Original references}
\makeatletter
\newlabel{lemma-normal-closure-inside-normal}{{0BMG}{source}{}{sourcepointer-4}{}}
\newlabel{lemma-normal-closure-inside-normal@cref}{{[originalreference][0][]0BMG}{[1][0][]source}{}{}{}}
\makeatother
\crefname{originalreference}{original reference}{original references}
\Crefname{originalreference}{Original reference}{Original references}
\makeatletter
\newlabel{theorem-galois-theory}{{09DW}{source}{}{sourcepointer-5}{}}
\newlabel{theorem-galois-theory@cref}{{[originalreference][0][]09DW}{[1][0][]source}{}{}{}}
\makeatother
\title{Fundamental theorem of infinite Galois theory: closed subgroups correspond to subextensions}
\author{The Stacks Project Authors}
\date{stacks-project:fields.tex:local-snapshot-20261012; GNU FDL 1.2 or later}
\begin{document}
\maketitle
\noindent\textit{Editorial note.} All delivered mathematical text is the author's own, copied from the pinned Stacks Project chapter file fields.tex; no statement or proof step was rewritten, reordered or removed.\par\medskip
\noindent\textit{Editorial note.} The counted proof is the author's complete argument: The kernel of the restriction map Gal(L/K) -> Gal(E/K) is open because Gal(E/K) is discrete; the finite case E\textasciicircum{}\{Gal(E/K)\} = K is theorem-galois-theory of the same chapter.\par\medskip
\noindent\textit{Editorial note.} The label theorem-inifinite-galois-theory names the counted statement, and the proof environment immediately following it in the source is delivered verbatim as its proof.\par\medskip
\noindent\textit{Editorial note.} The counted statement is carried by the author's own theorem environment in the chapter source.\par\medskip
\noindent\textit{Allowed standard prerequisites (editorial).}\begin{itemize}
\item Same chapter theorem-galois-theory (finite fundamental theorem, delivered structure): finite Galois correspondence used throughout without further mention (author's announcement at source lines 3288-3290)
\item Same chapter lemma-galois-profinite (the profinite topology on Gal(L/K)), lemma-galois-infinite (canonical surjective continuous restriction map Gal(L/K) -> Gal(E/K)), lemma-galois-goes-up, lemma-normal-closure-inside-normal
\item Topology, Lemma topology-lemma-bijective-map is not used here; the closed subgroup argument uses only the topology definition and finite Galois theory
\end{itemize}\medskip
\ifdefined\hypertarget\hypertarget{sourcepointer-1}{}\fi
\noindent\textit{Source reference.} Original reference 09I2: sources/stacks-fields/pinned-original-source.tex, source lines 2911--2911. Prerequisite pointer invoked at 3305; it names a standard fact of the same chapter that the counted argument uses, and every step of the counted proof built on it is retained verbatim in the delivered body.\par\medskip
\ifdefined\hypertarget\hypertarget{sourcepointer-2}{}\fi
\noindent\textit{Source reference.} Original reference 0BMK: sources/stacks-fields/pinned-original-source.tex, source lines 3168--3168. Prerequisite pointer invoked at 3296; it names a standard fact of the same chapter that the counted argument uses, and every step of the counted proof built on it is retained verbatim in the delivered body.\par\medskip
\ifdefined\hypertarget\hypertarget{sourcepointer-3}{}\fi
\noindent\textit{Source reference.} Original reference 0BMJ: sources/stacks-fields/pinned-original-source.tex, source lines 3076--3076. Prerequisite pointer invoked at 3274; it names a standard fact of the same chapter that the counted argument uses, and every step of the counted proof built on it is retained verbatim in the delivered body.\par\medskip
\ifdefined\hypertarget\hypertarget{sourcepointer-4}{}\fi
\noindent\textit{Source reference.} Original reference 0BMG: sources/stacks-fields/pinned-original-source.tex, source lines 2351--2351. Prerequisite pointer invoked at 3293; it names a standard fact of the same chapter that the counted argument uses, and every step of the counted proof built on it is retained verbatim in the delivered body.\par\medskip
\ifdefined\hypertarget\hypertarget{sourcepointer-5}{}\fi
\noindent\textit{Source reference.} Original reference 09DW: sources/stacks-fields/pinned-original-source.tex, source lines 3001--3001. Prerequisite pointer invoked at 3289; it names a standard fact of the same chapter that the counted argument uses, and every step of the counted proof built on it is retained verbatim in the delivered body.\par\medskip
\begin{theorem}[Fundamental theorem of infinite Galois theory]
\label{theorem-inifinite-galois-theory}
Let $L/K$ be a Galois extension. Let $G = \text{Gal}(L/K)$
be the Galois group viewed as a profinite topological group
(Lemma \ref{lemma-galois-profinite}). Then we have $K = L^G$ and the map
$$
\{\text{closed subgroups of }G\}
\longrightarrow
\{\text{subextensions }L/M/K\},\quad
H \longmapsto L^H
$$
is a bijection whose inverse maps $M$ to $\text{Gal}(L/M)$.
The finite subextensions $M$ correspond exactly to the open
subgroups $H \subset G$. The normal closed subgroups $H$ of $G$
correspond exactly to subextensions $M$ Galois over $K$.
\end{theorem}

\begin{proof}
We will use the result of finite Galois theory
(Theorem \ref{theorem-galois-theory})
without further mention.
Let $S \subset L$ be a finite subset. There exists a tower
$L/E/K$ such that $K(S) \subset E$ and such that
$E/K$ is finite Galois, see Lemma \ref{lemma-normal-closure-inside-normal}.
In other words, we see that $L/K$ is the union of its finite
Galois subextensions.
For such an $E$, by Lemma \ref{lemma-galois-infinite}
the map $\text{Gal}(L/K) \to \text{Gal}(E/K)$ is surjective
and continuous, i.e., the kernel is open because the topology
on $\text{Gal}(E/K)$ is discrete.
In particular we see that no element of $L \setminus K$ is fixed by
$\text{Gal}(L/K)$ as $E^{\text{Gal}(E/K)} = K$.
This proves that $L^G = K$.

\medskip\noindent
By Lemma \ref{lemma-galois-goes-up} given a subextension $L/M/K$
the extension $L/M$ is Galois. It is immediate from the definition
of the topology on $G$ that the subgroup $\text{Gal}(L/M)$ is closed.
By the above applied to $L/M$ we see that $L^{\text{Gal}(L/M)} = M$.

\medskip\noindent
Conversely, let $H \subset G$ be a closed subgroup. We claim that
$H = \text{Gal}(L/L^H)$. The inclusion $H \subset \text{Gal}(L/L^H)$
is clear. Suppose that $g \in \text{Gal}(L/L^H)$. Let $S \subset L$
be a finite subset. We will show that the open neighbourhood
$U_S(g) = \{g' \in G \mid g'(s) = g(s)\}$ of $g$ meets $H$.
This implies that $g \in H$ because $H$ is closed.
Let $L/E/K$ be a finite Galois subextension containing $K(S)$
as in the first paragraph of the proof and consider the homomorphism
$c : \text{Gal}(L/K) \to \text{Gal}(E/K)$.
Then $L^H \cap E = E^{c(H)}$. Since $g$ fixes $L^H$ it fixes
$E^{c(H)}$ and hence $c(g) \in c(H)$ by finite Galois theory.
Pick $h \in H$ with $c(h) = c(g)$. Then $h \in U_S(g)$ as desired.

\medskip\noindent
At this point we have established the correspondence between closed
subgroups and subextensions.

\medskip\noindent
Assume $H \subset G$ is open. Arguing as above we find that
$H$ contains $\text{Gal}(L/E)$ for some large enough finite
Galois subextension $E$ and we find that $L^H$ is contained
in $E$ whence finite over $K$. Conversely, if $M$ is a finite
subextension, then $M$ is generated by a finite subset $S$
and the corresponding subgroup is the open subset $U_S(e)$
where $e \in G$ is the neutral element.

\medskip\noindent
Assume that $K \subset M \subset L$ with $M/K$ Galois.
By Lemma \ref{lemma-galois-infinite} there is a surjective
continuous homomorphism of Galois groups
$\text{Gal}(L/K) \to \text{Gal}(M/K)$ whose
kernel is $\text{Gal}(L/M)$. Thus $\text{Gal}(L/M)$ is a normal
closed subgroup.

\medskip\noindent
Finally, assume $N \subset G$ is normal and closed. For any
$L/E/K$ as in the first paragraph of the proof, the image
$c(N) \subset \text{Gal}(E/K)$ is a normal subgroup.
Hence $L^N = \bigcup E^{c(N)}$ is a union of Galois extensions
of $K$ (by finite Galois theory) whence Galois over $K$.
\end{proof}

\end{document}
